\ifdefined\mainfile\else
\documentclass{article}
\usepackage{graphicx}
\usepackage{float}
\usepackage{amsmath}
\usepackage{amssymb}
\usepackage[a4paper, left=2.5cm, right=2.5cm, top=2.5cm, bottom=2.5cm]{geometry}
\usepackage[colorlinks=true,linkcolor=black,citecolor=blue,urlcolor=blue]{hyperref}
\begin{document}
\fi

\section{Software design}
\label{sec:software}

The software has three layers: firmware on the ESP8266, the ThingSpeak cloud channel that mediates between everyone, and the static web dashboard hosted on Netlify. Each layer is independent -- if any one is restarted, the others keep working.

\subsection{Firmware architecture (ESP8266)}
\label{sec:firmware}

The firmware is written in Arduino C++ and built around an explicit state machine. The main loop dispatches on a \texttt{State} variable and the four normal states (\textsc{load}, \textsc{check}, \textsc{close}, \textsc{empty\_me}) are joined by an error state (\textsc{help}) for overfill and stuck-motor conditions. Three independent concerns run alongside the state machine:

\begin{itemize}
    \item \textbf{Sensor sampling.} The HC-SR04 is read by sending a 10\,$\mu$s trigger pulse and timing the echo with \texttt{pulseIn()}; distance is converted using the speed of sound ($d = t \cdot 0.0343 / 2$ cm). The LM35 is sampled through the ESP's 10-bit ADC with the reference calibrated to 2.667\,V, and the raw reading is smoothed with an IIR filter (10-sample running mean blended 30/70 with the previous value).
    \item \textbf{ThingSpeak push.} A non-blocking timer pushes fields 1--4 and 6 every 16\,s when data has changed (the free-tier minimum is 15\,s between writes).
    \item \textbf{ThingSpeak poll.} A separate timer polls field 5 (the bag-count command from the dashboard) every 30\,s, leaving CPU headroom for the state machine.
\end{itemize}

To save power between pushes the WiFi radio is dropped into light-sleep, which cuts idle current significantly while still allowing the ESP to wake on a timer.

\subsection{Cloud layer (ThingSpeak)}
\label{sec:cloud}

ThingSpeak channel 3364403 acts as the single source of truth. Six of its eight fields are in use; the mapping is shown in Table~\ref{tab:fields}.

\begin{table}[H]
    \centering
    \begin{tabular}{c|l|l|l}
        \textbf{Field} & \textbf{Meaning} & \textbf{Written by} & \textbf{Read by} \\
        \hline
        1 & Distance (cm) from ultrasonic sensor & ESP & Dashboard \\
        2 & Bags remaining                       & ESP & Dashboard, IFTTT React \\
        3 & State (0=LOAD, 1=CHECK, 2=CLOSE, 3=EMPTY\_ME) & ESP & Dashboard \\
        4 & Bag-full event (1 = full)            & ESP & IFTTT React \\
        5 & Set-bag-count command                & Dashboard & ESP \\
        6 & Temperature (\textdegree C) from LM35 & ESP & Dashboard \\
    \end{tabular}
    \caption{ThingSpeak field map}
    \label{tab:fields}
\end{table}

\subsection{Dashboard (static HTML on Netlify)}
\label{sec:dashboard}

The dashboard at \url{https://dtutrash.netlify.app/} is a single HTML file (\texttt{docs/index.html}) of plain HTML/CSS/JavaScript -- no framework, no backend, no build step. Netlify just serves the file. The page calls the ThingSpeak feeds endpoint on a 5\,s timer, picks the most recent entry that actually contains ESP data (ignoring entries that only set field 5) and updates the cards. A POST-style write goes back the other way when the user clicks Save on the bag-count input.

\subsection{Code, libraries and source}
\label{sec:code}

All firmware and dashboard code is in the project repository on GitHub. Third-party libraries used:

\begin{itemize}
    \item \texttt{ESP8266WiFi}, \texttt{ESP8266HTTPClient}, \texttt{ESP8266WebServer} -- WiFi and HTTP client (built into the ESP8266 Arduino core).
    \item \texttt{Stepper} (Arduino built-in) -- driving the 28BYJ-48 motors.
    \item \texttt{Adafruit\_GFX} + \texttt{Adafruit\_SSD1306} -- OLED rendering.
    \item ThingSpeak HTTP REST API -- no client library; we issue raw \texttt{http.GET()} calls.
\end{itemize}

The submitted DTU LEARN bundle includes the complete Arduino sketch (\texttt{sketch\_apr16a.ino}) and the dashboard source (\texttt{docs/index.html}), as well as IFTTT/ThingSpeak React configuration screenshots so the cloud-side configuration is reproducible.

\ifdefined\mainfile\else
\end{document}
\fi
